% ============================================================
%  示例：知行读书 Agent 编排流水线 (Flowchart v4)
%  v4 改进：
%    - 字体缩小（footnotesize 主体）
%    - 总线汇流法（所有分支汇聚到一条水平线，整齐统一）
%    - 芯片简化为单行
%    - 节点宽度收窄
%
%  编译命令：xelatex example-zhixing-flowchart.tex
% ============================================================

\documentclass[tikz,border=14pt]{standalone}

\usepackage{fontspec}
\usepackage{ctex}
\usepackage{fontawesome5}
\usetikzlibrary{positioning, fit, backgrounds, arrows.meta, calc}
\usetikzlibrary{shapes.geometric}

% ============================================================
% 字体
% ============================================================
\setmainfont{Arial}
\setsansfont{Arial}
\setmonofont{Consolas}[Scale=0.85]
\setCJKsansfont{Microsoft YaHei}
\setCJKmainfont{Microsoft YaHei}

% ============================================================
% 配色（蓝色主题）
% ============================================================
\definecolor{PaperBorder}{HTML}{B0C4DE}
\definecolor{PaperFill}{HTML}{FFFFFF}
\definecolor{PaperTitle}{HTML}{2C3E50}
\definecolor{PaperMuted}{HTML}{7F8C8D}
\definecolor{AccentOrange}{HTML}{D35400}
\definecolor{PaperLayerBg}{HTML}{F5F8FA}
\definecolor{PaperLayerBorder}{HTML}{D8E3EE}
\definecolor{BranchChipFill}{HTML}{EFF6FF}
\definecolor{BranchChipBdr}{HTML}{BFDBFE}

% ============================================================
% 尺寸常量
% ============================================================
\def\NodeW{3.8cm}
\def\NodeH{0.8cm}
\def\ChipW{2.4cm}
\def\ChipH{0.55cm}

\tikzset{
  % ---- 开始/结束 ----
  startend/.style = {
    rectangle,
    draw = AccentOrange,
    fill = AccentOrange!8,
    line width = 0.9pt,
    rounded corners = 14pt,
    minimum width = 3.0cm,
    minimum height = 0.75cm,
    align = center,
    font = \footnotesize\sffamily\bfseries,
    text = PaperTitle,
    inner sep = 4pt,
  },
  % ---- 过程 ----
  process/.style = {
    rectangle,
    draw = PaperBorder,
    fill = PaperFill,
    line width = 0.8pt,
    rounded corners = 4pt,
    minimum width = \NodeW,
    minimum height = \NodeH,
    align = center,
    font = \footnotesize\sffamily,
    text = PaperTitle,
    inner sep = 5pt,
  },
  % ---- 决策 ----
  decision/.style = {
    diamond,
    draw = PaperBorder,
    fill = PaperLayerBg,
    line width = 0.8pt,
    aspect = 2.5,
    minimum width = 1.8cm,
    minimum height = 0.7cm,
    align = center,
    font = \footnotesize\sffamily\bfseries,
    text = PaperTitle,
    inner sep = 1pt,
  },
  % ---- 分支芯片 ----
  chip/.style = {
    rectangle,
    draw = BranchChipBdr,
    fill = BranchChipFill,
    line width = 0.6pt,
    rounded corners = 3pt,
    minimum width = \ChipW,
    minimum height = \ChipH,
    align = center,
    font = \tiny\sffamily,
    text = PaperTitle,
    inner sep = 3pt,
  },
  % ---- 分支组背景框 ----
  branch-bg/.style = {
    draw = PaperLayerBorder,
    fill = PaperLayerBg,
    line width = 0.5pt,
    dashed,
    dash pattern = on 3pt off 2pt,
    rounded corners = 6pt,
    inner sep = 8pt,
  },
  % ---- 线条样式 ----
  bus/.style = {
    PaperBorder, line width = 0.8pt,
  },
  bus-merge/.style = {
    PaperMuted, line width = 0.7pt, dashed, dash pattern = on 3pt off 2pt,
  },
  arr/.style = {
    ->, >=Stealth, bus,
  },
  arr-merge/.style = {
    ->, >=Stealth, bus-merge,
  },
}

\begin{document}

\begin{tikzpicture}

  % ==========================================================
  %  坐标规划（紧凑）
  %  X: 芯片 4列，中心距 3.0cm
  %  Y: 主流程每步 1.2cm 间距
  % ==========================================================
  \def\CxOne{-4.5}
  \def\CxTwo{-1.5}
  \def\CxThr{1.5}
  \def\CxFou{4.5}

  % ---- 主流程节点 ----
  \node[startend] (start) at (0, 0) {\faPlay\ 用户输入消息};

  \node[process] (s1) at (0, -1.2)
    {\faComment\ \textbf{1.~意图分类}\ \ {\tiny\color{PaperMuted}1--5ms}};

  \node[decision] (d1) at (0, -2.6) {意图};

  % ===== 第一组芯片 =====
  \node[chip] (br1) at (\CxOne, -4.2) {knowledge\_query};
  \node[chip] (br2) at (\CxTwo, -4.2) {deep\_discussion};
  \node[chip] (br3) at (\CxThr, -4.2) {teaching\_practice};
  \node[chip] (br4) at (\CxFou, -4.2) {casual\_chat};

  \node[process] (s3) at (0, -5.8)
    {\faChartLine\ \textbf{3.~难度调整}\ \ {\tiny\color{PaperMuted}$<$1ms}};

  \node[decision] (d2) at (0, -7.2) {难度};

  % ===== 第二组芯片 =====
  \node[chip] (act1) at (\CxOne, -8.8) {increase\_bloom};
  \node[chip] (act2) at (\CxTwo, -8.8) {decrease\_bloom};
  \node[chip] (act3) at (\CxThr, -8.8) {mark\_mastered};
  \node[chip] (act4) at (\CxFou, -8.8) {maintain};

  \node[process] (s4) at (0, -10.4)
    {\faDatabase\ \textbf{4.~上下文构建}};

  \node[process] (s5) at (0, -11.6)
    {\faCogs\ \textbf{5.~系统提示组装}};

  \node[process] (s6) at (0, -12.8)
    {\faBolt\ \textbf{6.~流式响应}};

  \node[process, fill=PaperLayerBg, draw=PaperLayerBorder, dashed]
    (post) at (0, -14.0)
    {\faCaretDown\ 异步后处理};

  \node[startend] (end) at (0, -15.2) {\faStop\ 返回完整响应};

  % ==========================================================
  %  分支组背景框
  % ==========================================================
  \begin{scope}[on background layer]
    \node[branch-bg, fit=(br1)(br2)(br3)(br4)] (box1) {};
    \node[branch-bg, fit=(act1)(act2)(act3)(act4)] (box2) {};
  \end{scope}

  % ==========================================================
  %  连接线 — 总线汇流法
  %  分散：源 → 垂直线 → 水平总线 → 垂直线 → 各芯片
  %  合并：各芯片 → 垂直线 → 水平总线 → 垂直线 → 目标
  % ==========================================================

  % --- 主流程 ---
  \draw[arr] (start) -- (s1);
  \draw[arr] (s1) -- (d1);

  % --- 决策1 → 总线分散 → 4芯片 ---
  \def\BusA{-3.5}  % 分散总线 Y
  \draw[bus] (d1.south) -- (0, \BusA);           % 垂直引线
  \draw[bus] (\CxOne, \BusA) -- (\CxFou, \BusA);  % 水平总线
  \draw[arr] (\CxOne, \BusA) -- (br1.north);      % 4条垂直分支线
  \draw[arr] (\CxTwo, \BusA) -- (br2.north);
  \draw[arr] (\CxThr, \BusA) -- (br3.north);
  \draw[arr] (\CxFou, \BusA) -- (br4.north);

  % --- 4芯片 → 总线合并 → s3 ---
  \def\BusB{-5.0}  % 合并总线 Y
  \draw[bus-merge] (br1.south) -- (\CxOne, \BusB);
  \draw[bus-merge] (br2.south) -- (\CxTwo, \BusB);
  \draw[bus-merge] (br3.south) -- (\CxThr, \BusB);
  \draw[bus-merge] (br4.south) -- (\CxFou, \BusB);
  \draw[bus-merge] (\CxOne, \BusB) -- (\CxFou, \BusB);  % 水平合并总线
  \draw[arr-merge] (0, \BusB) -- (s3.north);             % 统一出线

  \draw[arr] (s3) -- (d2);

  % --- 决策2 → 总线分散 → 4芯片 ---
  \def\BusC{-8.1}
  \draw[bus] (d2.south) -- (0, \BusC);
  \draw[bus] (\CxOne, \BusC) -- (\CxFou, \BusC);
  \draw[arr] (\CxOne, \BusC) -- (act1.north);
  \draw[arr] (\CxTwo, \BusC) -- (act2.north);
  \draw[arr] (\CxThr, \BusC) -- (act3.north);
  \draw[arr] (\CxFou, \BusC) -- (act4.north);

  % --- 4芯片 → 总线合并 → s4 ---
  \def\BusD{-9.6}
  \draw[bus-merge] (act1.south) -- (\CxOne, \BusD);
  \draw[bus-merge] (act2.south) -- (\CxTwo, \BusD);
  \draw[bus-merge] (act3.south) -- (\CxThr, \BusD);
  \draw[bus-merge] (act4.south) -- (\CxFou, \BusD);
  \draw[bus-merge] (\CxOne, \BusD) -- (\CxFou, \BusD);
  \draw[arr-merge] (0, \BusD) -- (s4.north);

  % --- 后续主流程 ---
  \draw[arr] (s4) -- (s5);
  \draw[arr] (s5) -- (s6);
  \draw[arr] (s6) -- (post);
  \draw[arr] (post) -- (end);

  % ==========================================================
  %  图例（右下角，简洁）
  % ==========================================================
  \node[font=\tiny\sffamily, text=PaperTitle, anchor=west]
    at ($(end.south east)+(0.5cm,-0.6cm)$) {
    \begin{tabular}{@{}r@{\ }l@{\quad}r@{\ }l@{}}
      {\color{PaperBorder}\rule[1pt]{10pt}{0.8pt}$>$} & 主流程
    & {\color{PaperMuted}\rule[1pt]{2pt}{0.8pt}\,{-}\,\rule[1pt]{2pt}{0.8pt}\,{-}\,\rule[1pt]{2pt}{0.8pt}$>$} & 分支合并
    \end{tabular}
  };

\end{tikzpicture}

\end{document}
